\section{Invocation Modes and Execution Lifecycle}
\subsection{Discovery, Known Providers, and Local Composition}
Normal discovery is appropriate when the executor has not been chosen. The User chooses among positive ACKs that pass validation and policy checks. Selection timing follows the configured policy rather than one universal waiting rule.

A supported first-positive-ACK path may select before timeout. A collection-based policy waits until its declared closure condition is met. The proposed snapshot-based collaboration commits only after its candidate collection closes.

A negative ACK communicates non-participation; silence cannot distinguish refusal from loss or disconnection. Selecting several independent executions does not assign complementary roles and dependencies as a collaboration plan does.

Known-Provider invocation is appropriate when the operation must use a particular device, for example, stopping the camera that was previously started. The current \texttt{RequestServiceTargeted} path can reuse previously bootstrapped one-time authorization material to avoid repeated ACK/Selection exchanges under its supported conditions. It still checks current authorization and request state. The fast path is not universal: request-scoped confidentiality needs a Selection envelope to carry the input and response keys, and the current implementation uses the bootstrap/ACK/Selection path for that case. The explicitly Selection-gated input mode is rejected by the Selection-free Targeted interface. Thus, targeting a known Provider is an application requirement, whereas omitting exchanges is a conditional implementation optimization.

Trusted local composition operates within a process through the local registry. It has neither remote discovery nor the same network threat boundary. Evaluation should identify whether a handler was reached locally, by normal discovery, or by a targeted exchange. Otherwise an apparent invocation improvement might merely measure a different operation.

\subsection{Admission, Progress, and Failure}
An API submission, admission to a runtime queue, acceptance by a Provider, and completion of application work are different events. Under overload, reporting latency only for completed requests can hide dropped or rejected work. The runtime therefore needs observable outcomes for refusal, timeout, and terminal results, while an evaluation must count all submitted requests in its stated denominator. Selection-status queries can help determine the observed status of selected work; they are not proof of a physical effect or a distributed transaction.

Applications must also distinguish cancellation of interest in a result from reversal of an action. A recording can be stopped, but previously recorded frames still exist. An inference worker can stop producing new output, but a completed remote action cannot be undone by deleting requester state. Retrying requires the application's stated idempotency or compensation policy. For the proposed collaboration extension, replacement additionally requires old assignments and late callbacks to stop advancing the current plan.

\section{Control Messages, Streams, and Complete Objects}
Small service messages communicate intent, permission evidence, and status. Large application data should not be copied into each control message. An object reference allows a service to identify a model fragment, recording, or image, while the recipient retrieves its named Data separately. The reference and service contract must supply enough information to select the intended version and validate the reconstructed object. The concrete interface determines the available name, size, digest, segmentation, and protection metadata; not every object type needs an identical collection of fields.

The important separation is between identifying an object, obtaining its bytes, and authorizing its use. A digest can detect a mismatch with the expected bytes only when the expected digest itself is trusted. A signature binds Data to a key under the application's trust rules. The invocation or collaboration checks establish whether that producer's object is permitted as an input to the current execution. Recipient-scoped key delivery controls access to protected Content. None of these checks should silently substitute for another.

An ongoing stream has a different lifecycle from a complete object. Live video may prefer recent frames and tolerate explicitly identified gaps. A model artifact generally requires every needed byte before loading. Incremental inference requires ordered application events and a terminal outcome, but its token stream is not necessarily interchangeable with a camera stream interface. NDNSF supplies named-data mechanisms; applications define acceptable gaps, buffering limits, and completeness rules. The evaluation must measure the chosen policy, including its discarded data and queueing costs.

Named-object reuse is valuable in the multi-view case: the compute Provider returns annotations and references to the original images, and provides the requester with the image content keys it is permitted to receive. The requester can retrieve those images from available producers or caches and validate their original signatures. This avoids producing a second image object merely to change the receiver. It does not guarantee that a cached copy remains available; retention requirements need an explicit producer or repository policy.
